
\begin{abstract}
In \textbf{DiT-based video generation models equipped with 3D Rotary Position Embeddings (3D RoPE)}, the attention mechanism remains a primary computational bottleneck due to its quadratic complexity with respect to sequence length. While quantized \textbf{FlashAttention} offers a promising path toward hardware acceleration, existing low-bit quantization methods overlook two critical challenges in this setting: \textbf{1)} applying online rotation matrices---a widely used technique for mitigating outliers in Queries ($Q$) and Keys ($K$)---is difficult to reconcile with \textbf{RoPE}; and \textbf{2)} the non-negative attention matrix $P = \exp(QK - \max(QK))$ makes symmetric quantization waste half of the 4-bit dynamic range. In this work, we observe that the outlier distributions of $Q$ and $K$ are strongly affected by the dimensional partitioning of \textbf{3D RoPE}. Based on this finding, we propose \textbf{RotateAttention}, an efficient \textbf{mixed-precision INT4 FlashAttention} framework tailored for \textbf{DiT-based video generation models with 3D RoPE}, using selective \textbf{FP16 fallback} for accuracy-sensitive attention blocks and denoising steps. RotateAttention introduces two core techniques: \textbf{1) RoPE-aware Rotation}, which employs either mergeable rotation matrices that can be fused into RoPE or negligible-overhead matrices to mitigate RoPE-induced outliers in $Q$ and $K$; and \textbf{2) Range-optimized $P$ Quantization}, which uses fixed scales and zero-points to fully exploit the \textbf{INT4 numerical range} with minimal computational overhead. 
Experiments show that \textbf{RotateAttention} preserves video generation quality nearly identical to full-precision baselines while achieving up to 1.68$\times$ end-to-end speedup and 2.2$\times$ kernel-level acceleration.

\keywords{Video generation \and Quantized attention \and 3D RoPE}
\end{abstract}


